\documentclass[11pt,a4paper]{article}

% A realistic single-file paper for the v0.1 exit demo (DESIGN.md §7):
% open it, edit a sentence, watch the PDF update. It uses only packages Tectonic fetches on
% its own, so it also proves "zero setup to first PDF".

\usepackage[T1]{fontenc}
\usepackage{lmodern}
\usepackage{amsmath,amssymb,amsthm}
\usepackage{booktabs}
\usepackage{microtype}
\usepackage[hidelinks]{hyperref}

\newtheorem{theorem}{Theorem}
\newtheorem{lemma}[theorem]{Lemma}
\theoremstyle{definition}
\newtheorem{definition}[theorem]{Definition}

\title{Local Compilation as a Design Principle for Scientific Writing Tools}
\author{A.~Author \and B.~Coauthor}
\date{September 2026}

\begin{document}
\maketitle

\begin{abstract}
We argue that the dominant architecture for collaborative \LaTeX{} editing, in which a shared
server compiles every document, imposes a latency floor that no client-side optimisation can
remove. Moving compilation to the author's machine removes the floor entirely. We describe the
consequences of that move for the remaining components of an editor, quantify the latency budget
that becomes available, and outline a document model that keeps plain files as the source of
truth while still admitting real-time collaboration.
\end{abstract}

\section{Introduction}

Every editing environment for \LaTeX{} makes one decision before all others: where the
typesetting engine runs. Hosted services run it on a shared pool of workers behind a queue.
Desktop editors run it on the author's machine. The two choices look interchangeable from the
outside, since both produce a PDF, but they differ in a way that determines nearly everything
about the resulting product.

The queue is the problem. A shared compiler must schedule jobs, and scheduling means waiting.
The wait is not an inefficiency that can be optimised away; it is the cost of sharing. An editor
that compiles locally does not pay it, and an editor that does not pay it can spend the same
budget on the writing experience instead.

This paper makes three contributions. First, we decompose the compile latency of a hosted
service into stages and show which stages vanish under local compilation
(Section~\ref{sec:latency}). Second, we describe the three subsystems that a local editor
should invest the recovered budget in (Section~\ref{sec:editor}). Third, we sketch a document
model that reconciles two apparently contradictory requirements: that the file on disk remains a
plain text file, and that the session model supports concurrent editing
(Section~\ref{sec:model}).

\section{The latency budget}
\label{sec:latency}

Let a compile cycle consist of $n$ stages with durations $t_1, \dots, t_n$. The author perceives
the total
\begin{equation}
  T = \sum_{i=1}^{n} t_i ,
  \label{eq:total}
\end{equation}
and only $T$ matters: a stage that is fast in isolation is worthless if the stage before it is a
queue. Table~\ref{tab:stages} lists the stages for a typical hosted service and for a local
engine with a warm cache.

\begin{table}[ht]
  \centering
  \caption{Illustrative compile-cycle stages. Bold rows exist only under sharing.}
  \label{tab:stages}
  \begin{tabular}{lrr}
    \toprule
    Stage & Hosted (s) & Local, warm (s) \\
    \midrule
    Sync or debounce        & 0.3 & 0.7 \\
    \textbf{Queue wait}     & \textbf{1.5} & --- \\
    \textbf{Container start} & \textbf{2.0} & --- \\
    Compile                 & 6.0 & 0.8 \\
    Transfer and render     & 0.7 & 0.06 \\
    \midrule
    Total                   & 10.5 & 1.6 \\
    \bottomrule
  \end{tabular}
\end{table}

The bold rows are the argument. They are not slow implementations of necessary work; they are
work that exists because a compiler is shared. Equation~\eqref{eq:total} is indifferent to why a
stage exists, and so is the author.

\subsection{What a warm cache buys}

A local engine keeps its auxiliary files between runs. Cross-reference resolution, bibliography
processing and font loading all amortise across edits, and a precompiled preamble removes the
cost of reading package files on every cycle. Let $c$ be the cold compile time and $w$ the warm
one; empirically $w / c \approx 0.15$ for a document with a heavy preamble, and the ratio
improves as the preamble grows.

\begin{definition}[Incremental cycle]
  A compile cycle is \emph{incremental} if no stage repeats work whose inputs are unchanged
  since the previous cycle.
\end{definition}

\begin{theorem}
  Under local compilation with a warm cache, the perceived latency of an incremental cycle is
  bounded below by the debounce interval and above by the debounce interval plus the warm
  compile time.
\end{theorem}

\begin{proof}
  The debounce interval $d$ is deliberate and precedes every cycle, so $T \geq d$. With no queue
  and no container start, the remaining stages are the compile itself and the render, and the
  render of a document that fits in memory is negligible against the compile. Hence
  $T \leq d + w + \varepsilon$ for a small render cost $\varepsilon$.
\end{proof}

\section{Where the recovered budget goes}
\label{sec:editor}

Once latency is no longer the constraint, the remaining complaint about \LaTeX{} editors is that
they are thin: a text area with syntax colouring. Three deficits account for most of the
friction serious authors report.

\begin{enumerate}
  \item \textbf{Errors are shown, not explained.} The engine's log format predates the authors
    who read it. A message such as \texttt{Missing \$ inserted} is true and unhelpful.
  \item \textbf{Citations live elsewhere.} Reference management happens in a separate
    application, a separate file, and a publisher's web page, with the cost measured in
    window switches per paragraph.
  \item \textbf{Prose gets no help.} Autocompletion for commands is not assistance with the
    sentence.
\end{enumerate}

Each deficit maps to a subsystem: a diagnostic translator with a catalogue of rules, a
bibliography layer that acquires entries from identifiers, and a selection-scoped prose
assistant that proposes diffs rather than emitting text. We treat the first as the
differentiator: it is deterministic, works offline, and every rule added moves an error out of
the long tail permanently.

\section{A document model with plain files}
\label{sec:model}

Two requirements pull in opposite directions. The file on disk must remain plain \LaTeX{},
readable by any tool and by any collaborator who does not use ours. The session model must be a
conflict-free replicated data type, or real-time collaboration will require a rewrite later.

The resolution is to make the CRDT a \emph{session} model and never a storage format. Disk is
authoritative. Writes are whole-file and debounced. External changes, such as a version-control
checkout, are diffed and applied as CRDT transactions so that undo history and remote peers
remain coherent. When the file changes on disk while the buffer holds unsaved edits, the loop
stops and the author is asked; nothing about a manuscript is merged automatically.

\begin{lemma}
  If every write is whole-file and every external change is applied as a transaction, then at
  every quiescent moment the CRDT text equals the file content.
\end{lemma}

\begin{proof}[Proof sketch]
  Induct on the sequence of events. A local write copies CRDT text to disk; an external change
  copies disk text into the CRDT via a diff whose application is exact. Neither event can leave
  the two unequal once it completes, and quiescence means no event is in flight.
\end{proof}

\section{Related work}

Hosted collaborative editors popularised real-time \LaTeX{} editing and demonstrated demand for
it~\cite{hosted}. Desktop editors have compiled locally for decades but treated collaboration as
out of scope~\cite{desktop}. Operational transformation and CRDTs supply the theory of
concurrent editing~\cite{crdt}; our contribution is the constraint that the CRDT never becomes
the file.

\section{Conclusion}

Compilation belongs on the author's machine. That single choice removes the queue, and with it
the latency floor that hosted editors cannot lower. The editor that results should spend its
budget on explaining errors, acquiring citations and helping with prose, and it should do so
without ever taking the manuscript hostage in a format only it can read.

\begin{thebibliography}{9}
  \bibitem{hosted} A.~Vendor, ``A hosted collaborative editor for \LaTeX{},'' \emph{Proc.\ of
    Something}, 2014.
  \bibitem{desktop} B.~Maintainer, ``Twenty years of a desktop \TeX{} editor,'' \emph{TUGboat},
    vol.~40, no.~2, 2019.
  \bibitem{crdt} M.~Shapiro, N.~Pregui\c{c}a, C.~Baquero, and M.~Zawirski, ``Conflict-free
    replicated data types,'' in \emph{Stabilization, Safety, and Security of Distributed
    Systems}, 2011.
\end{thebibliography}

\end{document}
